This section tests how sensitive the results are to the main design choices. Each setting was trained with seeds 0 and 1 on the four satellite links, with all other settings as in Section~\ref{sec:protocol}, and evaluated on the same test sessions (80 additional runs, \texttt{results\_\allowbreak{}y3/\allowbreak{}sweeps/}, summarized by \texttt{experiments/\allowbreak{}y3\_\allowbreak{}sweeps.py}). The settings are a shorter (15 rows, 3\,s) and a longer (60 rows, 12\,s) input window, a longer horizon (30 steps, 6\,s), a smaller (64) and a larger (256) hidden size, and a pinball-loss weight of 0.25 and 4 instead of 1 in~\eqref{eq:loss}.

\subsection{Does the Year 3 gain depend on the setting?}
Table~\ref{tab:sweepvs} compares the Year 3 and Year 2 models trained with the same window and horizon. The interval-score gain holds in every setting and on every link, from 7\% to 41\%. The MAE gain also holds everywhere, and it is larger on OneWeb UL with the shorter or longer window than with the default. The RMSE changes by $-7.8\%$ to $+2.1\%$, with the largest reductions on OneWeb UL. The conclusions of Section~\ref{sec:results} therefore do not depend on the window or horizon chosen in Year 2.

\begin{table}[h]
\centering
\caption{Change (\%) of the Year 3 model against the Year 2 model trained with the same setting (RTT, seed means over two seeds, averaged over all horizon steps of the setting). Negative is better.}
\label{tab:sweepvs}
\scriptsize
\setlength{\tabcolsep}{2.5pt}
\begin{tabular}{l*{12}{r}}
\toprule
 & \multicolumn{3}{c}{OneWeb DL} & \multicolumn{3}{c}{OneWeb UL} & \multicolumn{3}{c}{Starlink DL} & \multicolumn{3}{c}{Starlink UL} \\
\cmidrule(lr){2-4}\cmidrule(lr){5-7}\cmidrule(lr){8-10}\cmidrule(lr){11-13}
Setting & MAE & RMSE & MIS & MAE & RMSE & MIS & MAE & RMSE & MIS & MAE & RMSE & MIS \\
\midrule
\tabinput{data/tab3_sweeps_vs.tex}
\bottomrule
\end{tabular}
\end{table}

\subsection{How the Year 3 model responds to each setting}
Table~\ref{tab:sweepds} gives the change of the Year 3 model under each setting relative to its default (first row, absolute values in ms). For the 30-step horizon, only the first 15 steps are scored. The numbers therefore compare the same 0.2--3\,s targets.
\begin{itemize}[leftmargin=*]
  \item \textbf{Input window.} A 12\,s window helps OneWeb UL (MAE $-9.4\%$) and hurts OneWeb DL (MAE $+11.2\%$). The 6\,s default is a reasonable compromise, and a per-link window is an easy gain for deployment.
  \item \textbf{Horizon.} Training for 6\,s costs little at 0.2--3\,s on three links (within about $\pm5\%$) but raises the Starlink DL MAE by 7.8\%. Figure~\ref{fig:sweepH} shows that the errors grow smoothly to 6\,s and that the Year 3 interval score stays below the Year 2 one at every step.
  \item \textbf{Model size.} A hidden size of 64 is worse on the downlinks (up to $+7.9\%$ MAE). A size of 256 changes the results by a few percent in both directions. The default of 128 is adequate.
  \item \textbf{Pinball weight.} A weight of 4 on the quantile loss improves the Year 3 model on every link and metric except the OneWeb DL RMSE ($+0.4\%$). The OneWeb DL MAE drops by 10.4\% and the interval score by 1.6\% to 5.9\%. A weight of 0.25 is worse on most links. The default weight of 1 used in the main study is therefore not the best choice, and the main results understate what the distributional model can reach. A weight of 4 is recommended for the next round.
\end{itemize}

\begin{table}[h]
\centering
\caption{Year 3 model under each setting: absolute RTT values for the default (first row, ms) and change (\%) against the default for the others (two seeds). The 6\,s horizon is scored on its first 15 steps.}
\label{tab:sweepds}
\scriptsize
\setlength{\tabcolsep}{2.5pt}
\begin{tabular}{l*{12}{r}}
\toprule
 & \multicolumn{3}{c}{OneWeb DL} & \multicolumn{3}{c}{OneWeb UL} & \multicolumn{3}{c}{Starlink DL} & \multicolumn{3}{c}{Starlink UL} \\
\cmidrule(lr){2-4}\cmidrule(lr){5-7}\cmidrule(lr){8-10}\cmidrule(lr){11-13}
Setting & MAE & RMSE & MIS & MAE & RMSE & MIS & MAE & RMSE & MIS & MAE & RMSE & MIS \\
\midrule
\tabinput{data/tab3_sweeps_ds2s.tex}
\bottomrule
\end{tabular}
\end{table}

\begin{figure}[h]
  \centering
  \resizebox{\linewidth}{!}{\begin{tikzpicture}
\begin{groupplot}[group style={group size=4 by 2, horizontal sep=1.2cm, vertical sep=1.6cm}, y3base, width=4.6cm, height=3.9cm, xmin=0.2, xmax=6, xtick={1,2,3,4,5,6}, xlabel={horizon (s)}]
\nextgroupplot[title={OneWeb DL}, ylabel={RTT MAE (ms)}, legend to name=legH, legend columns=2]
\addplot[mYtwo] table[x expr=\thisrow{h}*0.2, y=y2_seq2seq_mae] {data/sweepH_oneweb_dl.dat}; \addlegendentry{Y2 (2 seeds)}
\addplot[mYthree] table[x expr=\thisrow{h}*0.2, y=y3_dist_mae] {data/sweepH_oneweb_dl.dat}; \addlegendentry{Y3 (2 seeds)}
\nextgroupplot[title={OneWeb UL}]
\addplot[mYtwo] table[x expr=\thisrow{h}*0.2, y=y2_seq2seq_mae] {data/sweepH_oneweb_ul.dat};
\addplot[mYthree] table[x expr=\thisrow{h}*0.2, y=y3_dist_mae] {data/sweepH_oneweb_ul.dat};
\nextgroupplot[title={Starlink DL}]
\addplot[mYtwo] table[x expr=\thisrow{h}*0.2, y=y2_seq2seq_mae] {data/sweepH_starlink_dl.dat};
\addplot[mYthree] table[x expr=\thisrow{h}*0.2, y=y3_dist_mae] {data/sweepH_starlink_dl.dat};
\nextgroupplot[title={Starlink UL}]
\addplot[mYtwo] table[x expr=\thisrow{h}*0.2, y=y2_seq2seq_mae] {data/sweepH_starlink_ul.dat};
\addplot[mYthree] table[x expr=\thisrow{h}*0.2, y=y3_dist_mae] {data/sweepH_starlink_ul.dat};
\nextgroupplot[title={OneWeb DL}, ylabel={RTT interval score (ms)}]
\addplot[mYtwo] table[x expr=\thisrow{h}*0.2, y=y2_seq2seq_mis] {data/sweepH_oneweb_dl.dat};
\addplot[mYthree] table[x expr=\thisrow{h}*0.2, y=y3_dist_mis] {data/sweepH_oneweb_dl.dat};
\nextgroupplot[title={OneWeb UL}]
\addplot[mYtwo] table[x expr=\thisrow{h}*0.2, y=y2_seq2seq_mis] {data/sweepH_oneweb_ul.dat};
\addplot[mYthree] table[x expr=\thisrow{h}*0.2, y=y3_dist_mis] {data/sweepH_oneweb_ul.dat};
\nextgroupplot[title={Starlink DL}]
\addplot[mYtwo] table[x expr=\thisrow{h}*0.2, y=y2_seq2seq_mis] {data/sweepH_starlink_dl.dat};
\addplot[mYthree] table[x expr=\thisrow{h}*0.2, y=y3_dist_mis] {data/sweepH_starlink_dl.dat};
\nextgroupplot[title={Starlink UL}]
\addplot[mYtwo] table[x expr=\thisrow{h}*0.2, y=y2_seq2seq_mis] {data/sweepH_starlink_ul.dat};
\addplot[mYthree] table[x expr=\thisrow{h}*0.2, y=y3_dist_mis] {data/sweepH_starlink_ul.dat};
\end{groupplot}
\node at ($(group c2r2.south)!0.5!(group c3r2.south)+(0,-1.3cm)$) {\pgfplotslegendfromname{legH}};
\end{tikzpicture}
}
  \caption{Models trained for a 6\,s horizon: RTT MAE (top) and calibrated interval score (bottom) per step (seed means over two seeds).}
  \label{fig:sweepH}
\end{figure}
